We keep the notation of the previous subsection: real stability means nonvanishing when every variable lies in the open upper half-plane; $a_\lambda$ are the coefficients in the Schur expansion $p=\sum_{\lambda\vdash d}a_\lambda s_\lambda$ of a symmetric homogeneous $p$ of degree $d$; $f^\lambda$ counts standard Young tableaux of shape $\lambda$; and $\Imm_\lambda$ is the immanant of a $d\times d$ matrix attached to the irreducible character $\chi^\lambda$, with $\Imm_{(1^d)}=\det$ and $\Imm_{(d)}=\per$.  Problem 36 asks whether stability forces Schur-positivity outright; Problems 37 and 38 ask for the two endpoint inequalities that would bracket the Schur coefficients between the $\det$ and $\per$ ends of the immanant scale.
